% =====================================================================
% 04_partie_b.tex — manipulations de données
% =====================================================================
\chapter{Partie B : manipulations de données}
\label{chap:partieB}

\section{Objectifs et vue d'ensemble}

La partie B vérifie que le schéma de la partie A supporte des manipulations
réelles, y compris des cas limites. Elle ne se limite pas à des insertions
simples : elle démontre une contrainte différée, l'application d'une valeur par
défaut, la manipulation d'une \emph{nested table}, la mise à jour d'un
\code{VARRAY} et une suppression en cascade applicative. Six blocs, notés B5a
à B5f, structurent cette démonstration.

\begin{table}[htbp]
  \centering
  \small
  \caption{Les six manipulations de la partie B.}
  \label{tab:manipulations-B}
  \begin{tabularx}{\textwidth}{@{}C{1.0cm}L{4.7cm}L{4.5cm}X@{}}
    \toprule
    \entete{Bloc} & \entete{Action} & \entete{Technique démontrée} &
    \entete{Résultat attendu} \\
    \midrule
    B5a & Ajout du client 99 \code{OUKACI Ahmed} &
    Insertion d'un client avec un \code{VARRAY} d'un seul téléphone. &
    1 ligne créée. \\
    B5b & Création de \code{CC99} sans titulaire, puis \code{UPDATE} &
    Contrainte \code{REF IS NOT NULL} différée sur \code{cpt\_courant}. &
    Insertion acceptée, puis mise à jour, puis \code{COMMIT}. \\
    B5c & Opération de 50\,000~DA sans date &
    \code{DEFAULT SYSDATE} sur \code{date\_op}. &
    La date du jour est automatiquement mise. \\
    B5d & Signataire de droit \code{D} sur \code{CC99} &
    \code{INSERT INTO TABLE(SELECT signataires ...)}. &
    1 signataire ajouté dans la \emph{nested table}. \\
    B5e & \code{change\_Portable} &
    Mise à jour d'un élément de \code{VARRAY} avec gestion de \code{NULL},
    1, 2 ou 3 éléments. &
    Client 99 : 1 $\rightarrow$ 2 téléphones ; client 3 : 2\up{e} téléphone
    = \code{0777-11-22-33}. \\
    B5f & \code{drop\_client(99)} &
    Suppression en cascade applicative : opérations, signataires, filleuls,
    comptes, client. &
    Client 99 supprimé, aucun orphelin. \\
    \bottomrule
  \end{tabularx}
\end{table}

\begin{clebox}[Ce que B5b démontre exactement]
On peut insérer un compte courant sans titulaire \emph{à l'intérieur d'une
transaction}, puis le mettre à jour avant le \code{COMMIT}. La contrainte
\code{ck\_cc\_client} est différée et vérifiée à la validation. C'est un cas
typique de contrainte différée : la règle métier est respectée à la fin de la
transaction, mais pas nécessairement après chaque instruction.
\end{clebox}

\section{B5a : ajout du client 99}

Le client 99 est ajouté avec un \code{VARRAY} d'un seul élément. Ce choix
prépare B5e : la procédure \code{change\_Portable} devra étendre le tableau
avant d'écrire le deuxième téléphone. Les valeurs non précisées par le sujet
— prénom, date de naissance, numéro de sécurité sociale et adresse électronique
— sont choisies de façon cohérente.

\begin{lstlisting}[caption={B5a : insertion du client 99 avec un téléphone.},label={lst:b5a}]
INSERT INTO client (num_client, nom, prenom, date_naissance, num_secu,
                    adresse, telephones, email, date_entree,
                    ref_conseiller, ref_parrain)
VALUES (99, 'OUKACI', 'Ahmed', DATE '1999-01-01', '990101099',
        adresse_type(10, 'Rue de la Poste', 'Bejaia', '06000'),
        tel_vry_type('0661-00-00-01'),
        'ahmed.oukaci@mail.dz', TRUNC(SYSDATE),
        (SELECT REF(c) FROM conseiller c WHERE c.num_conseiller = 1), NULL);
COMMIT;
\end{lstlisting}

Le contrôle consiste à vérifier la ligne créée, puis à interroger le
\code{VARRAY} avec \code{TABLE(...)}. Cette seconde requête montre que le
\code{VARRAY} est bien un attribut multivalué et non une chaîne de caractères.

\begin{lstlisting}[caption={Contrôle du client 99 et de son téléphone.},label={lst:b5a-controle}]
SELECT c.num_client, c.nom, c.prenom, c.adresse.ville AS ville
  FROM client c
 WHERE c.num_client = 99;

SELECT t.COLUMN_VALUE AS telephone
  FROM TABLE (SELECT telephones FROM client WHERE num_client = 99) t;
\end{lstlisting}

\section{B5b : contrainte différée sur le titulaire}

B5b crée le compte \code{CC99} sans titulaire, puis le met à jour avant le
\code{COMMIT}. L'insertion serait refusée si la contrainte était immédiate.
Elle est acceptée parce que \code{ck\_cc\_client} est
\code{DEFERRABLE INITIALLY DEFERRED}. Oracle vérifie la contrainte au moment de
la validation, lorsque le titulaire est renseigné.

\begin{lstlisting}[caption={B5b : insertion sans titulaire, puis mise à jour.},label={lst:b5b}]
INSERT INTO cpt_courant (num_compte, solde, date_ouverture, ref_client,
                         ref_agence, decouvert, signataires)
VALUES ('CC99', 0, SYSDATE, NULL,
        (SELECT REF(a) FROM agence a WHERE a.code = 'A1'),
        10000, signataire_nt_type());

UPDATE cpt_courant
   SET ref_client = (SELECT REF(c) FROM client c WHERE c.num_client = 99)
 WHERE num_compte = 'CC99';

COMMIT;
\end{lstlisting}

\begin{attentionbox}[Ne pas oublier le \code{COMMIT}]
Tant que le \code{COMMIT} n'est pas exécuté, la contrainte différée n'est pas
vérifiée. Si l'on tente de valider sans avoir renseigné le titulaire, Oracle
lève une erreur d'intégrité. Le \code{ROLLBACK} annule l'insertion et remet la
base dans son état antérieur.
\end{attentionbox}

\section{B5c : valeur par défaut de la date}

B5c insère une opération de 50\,000~DA sans préciser \code{date\_op} ni
\code{commission}. L'insertion utilise une liste de colonnes, et non un
constructeur \code{operation\_type(...)}. Les valeurs par défaut déclarées sur
la table s'appliquent donc aux colonnes omises : \code{SYSDATE} pour la date et
0 pour la commission.

\begin{lstlisting}[caption={B5c : opération sans date explicite.},label={lst:b5c}]
INSERT INTO operation (id_op, ref_client, ref_compte, montant, type_op)
VALUES (seq_operation.NEXTVAL,
        (SELECT REF(c) FROM client c WHERE c.num_client = 99),
        (SELECT REF(x) FROM cpt_courant x WHERE x.num_compte = 'CC99'),
        50000, 'DEPOT');
COMMIT;

SELECT o.id_op, o.type_op, o.montant, o.commission,
       TO_CHAR(o.date_op, 'DD/MM/YYYY HH24:MI:SS') AS date_op,
       TO_CHAR(SYSDATE,  'DD/MM/YYYY HH24:MI:SS') AS date_sys
  FROM operation o
 WHERE o.ref_compte = (SELECT REF(x) FROM cpt_courant x
                       WHERE x.num_compte = 'CC99');
\end{lstlisting}

\begin{remarquebox}[Constructeur ou liste de colonnes ?]
Un constructeur \code{operation\_type(...)} oblige à fournir toutes les
valeurs, y compris la date. La liste de colonnes laisse Oracle appliquer les
valeurs par défaut. Pour démontrer \code{DEFAULT SYSDATE}, il faut donc
privilégier l'insertion par colonnes.
\end{remarquebox}

\section{B5d : signataire dans une \emph{nested table}}

L'ajout d'un signataire ne se fait pas par un \code{INSERT} classique sur une
table : la \emph{nested table} est un attribut de \code{cpt\_courant}. La
syntaxe \code{INSERT INTO TABLE (SELECT signataires ...)} permet d'insérer un
élément dans la collection d'une ligne précise.

\begin{lstlisting}[caption={B5d : ajout d'un signataire de droit \code{D}.},label={lst:b5d}]
INSERT INTO TABLE (SELECT signataires FROM cpt_courant
                    WHERE num_compte = 'CC99')
SELECT signataire_elt_nt_type(REF(c), 'D')
  FROM client c
 WHERE c.num_client = 1;
COMMIT;

SELECT s.ref_client.nom AS signataire, s.droit
  FROM cpt_courant c, TABLE(c.signataires) s
 WHERE c.num_compte = 'CC99';
\end{lstlisting}

\begin{attentionbox}[Pas de trigger sur la \emph{nested table}]
Oracle interdit de créer un trigger sur le stockage d'une \emph{nested table}.
Les contrôles de droits sont donc assurés par les méthodes
\code{estSignataire}, \code{peutRetirer} et par le trigger sur
\code{operation}. Cette limitation est assumée et documentée.
\end{attentionbox}

\section{B5e : mise à jour d'un \code{VARRAY}}

La procédure \code{change\_Portable} met à jour le deuxième téléphone d'un
client. Elle doit gérer quatre situations : \code{VARRAY} \code{NULL},
\code{VARRAY} vide, \code{VARRAY} à un seul élément et \code{VARRAY} à deux ou
trois éléments. Dans le cas à un élément, un \code{EXTEND(1)} est nécessaire
avant d'écrire à l'indice 2 ; sinon, Oracle lève \code{ORA-06533}.

\begin{lstlisting}[caption={Procédure \code{change\_Portable}.},label={lst:change-portable}]
CREATE OR REPLACE PROCEDURE change_Portable (
  paramcli IN NUMBER,
  paramtel IN VARCHAR2
) IS
  v_tel tel_vry_type;
BEGIN
  SELECT telephones INTO v_tel
    FROM client
   WHERE num_client = paramcli;

  IF v_tel IS NULL THEN
    v_tel := tel_vry_type(NULL, paramtel);
  ELSIF v_tel.COUNT = 0 THEN
    v_tel := tel_vry_type(NULL, paramtel);
  ELSIF v_tel.COUNT = 1 THEN
    v_tel.EXTEND(1);
    v_tel(2) := paramtel;
  ELSE
    v_tel(2) := paramtel;
  END IF;

  UPDATE client SET telephones = v_tel WHERE num_client = paramcli;
  COMMIT;

  DBMS_OUTPUT.PUT_LINE('change_Portable : client '||paramcli
                       ||' -> '||paramtel||'  (COUNT='||v_tel.COUNT||')');
EXCEPTION
  WHEN NO_DATA_FOUND THEN
    DBMS_OUTPUT.PUT_LINE('change_Portable : client '||paramcli
                         ||' inexistant.');
  WHEN OTHERS THEN
    ROLLBACK;
    DBMS_OUTPUT.PUT_LINE('change_Portable : echec, ROLLBACK ('
                         ||SQLERRM||')');
    RAISE;
END;
/
\end{lstlisting}

La procédure est autonome : elle valide en cas de succès et annule en cas
d'erreur. Ce choix est cohérent avec son usage, qui est une mise à jour
ponctuelle d'une donnée de contact. Il diffère de \code{virement}, qui ne
valide pas, car un virement doit pouvoir être enchaîné ou annulé par
l'appelant.

\begin{clebox}[Résultats attendus pour B5e]
Le client 99 passe d'un téléphone à deux téléphones : le premier reste
\code{0661-00-00-01}, le second devient \code{0770-99-88-77}. Le client 3,
qui possédait déjà deux entrées, voit son deuxième téléphone remplacé par
\code{0777-11-22-33}. Ces deux cas couvrent l'extension et le remplacement.
\end{clebox}

\section{B5f : suppression en cascade applicative}

La procédure \code{drop\_client} supprime un client et toutes ses dépendances.
Elle ne s'appuie pas sur un \code{ON DELETE CASCADE} déclaratif, mais sur une
séquence explicite de suppressions. Cet ordre est important : les opérations
référencent les comptes, les signataires référencent les clients, et les
filleuls référencent leur parrain.

\begin{enumerate}
  \item supprimer les opérations du client et celles de ses comptes ;
  \item retirer le client des listes de signataires des autres comptes ;
  \item mettre à \code{NULL} le \code{ref\_parrain} des clients qu'il
        parrainait ;
  \item supprimer ses comptes à terme, d'épargne et courants ;
  \item supprimer le client.
\end{enumerate}

\begin{lstlisting}[caption={Procédure \code{drop\_client}.},label={lst:drop-client}]
CREATE OR REPLACE PROCEDURE drop_client (paramcli IN NUMBER) IS
  v_ref REF client_type;
BEGIN
  SELECT REF(c) INTO v_ref FROM client c WHERE c.num_client = paramcli;

  DELETE FROM operation
   WHERE ref_client = v_ref
      OR ref_compte IN (SELECT REF(x) FROM cpt_courant x
                         WHERE x.ref_client = v_ref);

  FOR r IN (SELECT c.num_compte
              FROM cpt_courant c
             WHERE EXISTS (SELECT 1 FROM TABLE(c.signataires) s
                            WHERE s.ref_client = v_ref)) LOOP
    DELETE FROM TABLE (SELECT signataires FROM cpt_courant
                        WHERE num_compte = r.num_compte)
     WHERE ref_client = v_ref;
  END LOOP;

  UPDATE client SET ref_parrain = NULL WHERE ref_parrain = v_ref;

  DELETE FROM cpt_terme   WHERE ref_client = v_ref;
  DELETE FROM cpt_epargne WHERE ref_client = v_ref;
  DELETE FROM cpt_courant WHERE ref_client = v_ref;
  DELETE FROM client WHERE num_client = paramcli;

  COMMIT;
  DBMS_OUTPUT.PUT_LINE('drop_client : client '||paramcli
                       ||' supprime, COMMIT effectue.');
EXCEPTION
  WHEN NO_DATA_FOUND THEN
    DBMS_OUTPUT.PUT_LINE('drop_client : client '||paramcli||' inexistant.');
  WHEN OTHERS THEN
    ROLLBACK;
    DBMS_OUTPUT.PUT_LINE('drop_client : echec, ROLLBACK ('||SQLERRM||')');
    RAISE;
END;
/
\end{lstlisting}

\begin{attentionbox}[Pourquoi un \code{COMMIT} ici ?]
\code{drop\_client} est une opération de gestion autonome : elle doit persister.
\code{virement}, au contraire, est une brique réutilisable ; c'est l'appelant
qui décide de valider ou d'annuler. Les deux choix sont cohérents avec leur
usage. Cette distinction est une question classique de soutenance.
\end{attentionbox}

\section{Résultats et preuves}

Les résultats de la partie B sont visibles dans la sortie
\code{\_travail/\_B.out} et dans la section correspondante de
\code{\_travail/\_FULL.out}. Le tableau~\ref{tab:resultats-B} rappelle les
vérifications à effectuer.

\begin{table}[htbp]
  \centering
  \small
  \caption{Vérifications de la partie B.}
  \label{tab:resultats-B}
  \begin{tabularx}{\textwidth}{@{}C{1.2cm}L{6.4cm}X@{}}
    \toprule
    \entete{Bloc} & \entete{Requête de contrôle} & \entete{Résultat attendu} \\
    \midrule
    B5a & \code{SELECT * FROM client WHERE num\_client = 99;} &
    \code{OUKACI Ahmed}. \\
    B5b & \code{SELECT ... FROM cpt\_courant WHERE num\_compte = 'CC99';} &
    Compte sans titulaire, puis titulaire renseigné après \code{UPDATE}. \\
    B5c & \code{SELECT date\_op FROM operation WHERE ...;} &
    Date du jour. \\
    B5d & \code{SELECT ... FROM TABLE(SELECT signataires ...);} &
    Droit \code{D}. \\
    B5e & \code{SELECT ... FROM TABLE(c.telephones);} &
    Client 3 : \code{0777-11-22-33} en deuxième position. \\
    B5f & \code{SELECT COUNT(*) FROM client WHERE num\_client = 99;} &
    0. \\
    \bottomrule
  \end{tabularx}
\end{table}

\section{Pièges de la partie B}

Les erreurs les plus fréquentes de cette partie sont liées aux indices de
\code{VARRAY}, aux références dans les \emph{nested tables} et à l'ordre des
suppressions. Le tableau~\ref{tab:pieges-partieB} les récapitule.

\begin{table}[htbp]
  \centering
  \small
  \caption{Pièges de la partie B et solutions.}
  \label{tab:pieges-partieB}
  \begin{tabularx}{\textwidth}{@{}L{5.0cm}L{4.8cm}X@{}}
    \toprule
    \entete{Piège} & \entete{Symptôme} & \entete{Solution} \\
    \midrule
    \code{VARRAY} à un élément &
    \code{ORA-06533} sur l'indice 2. &
    Vérifier \code{COUNT}, puis \code{EXTEND}. \\
    Insertion d'un signataire &
    \code{ORA-22979} si le \code{REF} est basé sur la PK. &
    Garder l'OID système pour \code{client}. \\
    Suppression d'un client référencé &
    \code{ORA-02292} (FK). &
    Supprimer d'abord les dépendances ou les mettre à \code{NULL}. \\
    \code{COMMIT} au mauvais moment &
    Le test de rollback ne fonctionne plus. &
    Ne pas \code{COMMIT} dans \code{virement} ; \code{COMMIT} dans
    \code{drop\_client}. \\
    \bottomrule
  \end{tabularx}
\end{table}

\section{Conclusion de la partie B}

La partie B transforme un schéma statique en un système manipulable. Elle
démontre que les contraintes ne sont pas seulement décoratives : elles
autorisent des scénarios transactionnels précis, comme l'insertion différée
d'un compte sans titulaire. Elle montre aussi que les collections ne sont pas
des objets opaques : on peut les interroger, les étendre et les modifier. Enfin,
la suppression en cascade applicative prépare les questions de soutenance sur
l'intégrité référentielle et la différence entre \code{COMMIT} local et
transaction pilotée par l'appelant.
